
% ATLAS: Agentic Twin Ledger for Autonomous Systems
% George David Tsitlauri
% (add your affiliation/email if desired)
% Research conducted in 2026

\documentclass[conference]{IEEEtran}
\usepackage{graphicx}
\usepackage{amsmath}
\usepackage{hyperref}
\usepackage{listings}
\usepackage{booktabs}
\usepackage{multirow}

\title{ATLAS: Agentic Twin Ledger for Autonomous Systems\\\vspace{0.5em}\large{Research Year: 2026}}

\author{\IEEEauthorblockN{George David Tsitlauri}
\IEEEauthorblockA{\textit{Dept. of Informatics \& Telecommunications}\\
\textit{University of Thessaly, Greece}\\
gdtsitlauri@gmail.com}}

\begin{document}

\maketitle

\begin{abstract}
ATLAS is a research-grade, simulation-first prototype for decentralized cloud autonomy. It enables digital twins, intelligent agents, constrained LLM planning, RL adaptation, and trust-aware governance in a permissioned ledger, targeting safe, auditable, and decentralized control in distributed systems. This paper presents the architecture, methodology, experimental evaluation, and future directions of ATLAS.\\
\textbf{All research and experiments were conducted in 2026.}
\end{abstract}

\section{Introduction}
Modern distributed systems require autonomous control that is safe, auditable, and decentralized. ATLAS addresses these needs by combining digital twins, local agents, LLM-based planning, RL-guided adaptation, and a permissioned ledger for trusted coordination. The system is optimized for reproducible experiments and research.

\section{Related Work}
% Add references to digital twins, RL in distributed systems, LLM planning, decentralized governance, etc.
Recent advances in digital twins, reinforcement learning for cloud control, and decentralized ledgers have enabled new paradigms for autonomous systems. ATLAS builds on these by integrating LLM-based planning and trust-aware governance.

\section{System Architecture}
\subsection{Core Components}
\begin{itemize}
    \item \textbf{Cloud Simulator}: Models services, nodes, workload, faults, recovery, SLA.
    \item \textbf{Twin Runtime}: Local state mirror, history, what-if simulation.
    \item \textbf{Agent Core}: Proposal generation, peer validation, action interface.
    \item \textbf{LLM Planner}: Mock/Ollama/OpenAI-compatible planning with schema validation.
    \item \textbf{RL Engine}: Tabular Q-learning for adaptation.
    \item \textbf{Policy Guard}: Hard safety constraints and rejection logic.
    \item \textbf{Governance Chain}: Permissioned SQLite-backed append-only ledger.
    \item \textbf{Observability}: Metrics, events, traces, trust, RL stats, metadata.
    \item \textbf{API}: FastAPI runtime control-plane.
    \item \textbf{Dashboard}: Streamlit visibility layer.
\end{itemize}

\subsection{Interaction Diagrams}
% Insert mermaid/sequence diagrams as figures if desired

\section{Methodology}
\subsection{TRIDENT Algorithm}
TRIDENT (Trust-weighted Reinforcement-Informed Decentralized Twin Governance) computes action scores as:
\begin{align*}
\text{Score}(a) = &\; \alpha \cdot \text{TwinSimGain} + \beta \cdot \text{RLValue} + \gamma \cdot \text{SLAImprovement} \\
& - \lambda \cdot \text{Risk} - \mu \cdot \text{Cost} + \nu \cdot \text{Trust}
\end{align*}
Implemented in the agent core, this enables flexible, trust-aware decision-making.

\subsection{Baselines}
ATLAS supports multiple baselines: random policy, rule-based, ablated TRIDENT variants, and full TRIDENT.

\section{Experimental Evaluation}
\subsection{Setup}
Experiments are run on a Windows 11 laptop, using the provided scripts and configuration files. Benchmarks cover multiple scenarios (overload, node failure, latency spike, conflicting proposals, resource scarcity) and baselines.

\subsection{Metrics}
Collected metrics include decision/consensus latency, SLA violations, recovery time, resource utilization, cost, trust score evolution, RL reward trends, action success rate, and governance overhead.

\subsection{Results}
% Summarize key results, e.g., performance across baselines, effect of trust/RL, reproducibility, etc.
ATLAS demonstrates robust, reproducible control and governance in simulated distributed environments. The TRIDENT algorithm outperforms baselines in SLA adherence and trust evolution.

\section{Safety and Validation}
Safety constraints are enforced at multiple stages: proposal scoring, peer validation, and pre-execution. LLM outputs are schema-validated. Unsafe actions are rejected and logged.

\section{Reproducibility}
ATLAS ensures reproducibility via deterministic seeds, config snapshots, pinned dependencies, and versioned scenarios.

\section{Limitations and Future Work}
Current limitations include simulation-first design, lightweight governance, and compact RL. Future work includes richer trust models, adversarial simulations, distributed message-delay, deeper RL, and protocol variants.

\section{Conclusion}
ATLAS provides a modular, research-grade platform for decentralized cloud autonomy, integrating digital twins, RL, LLM planning, and trust-aware governance. It enables safe, auditable, and reproducible experimentation in distributed systems.

\section*{Acknowledgments}
% Add if needed

\bibliographystyle{IEEEtran}
\bibliography{atlas_refs}

\end{document}
